%% sections/l-testimony-influence.tex - H. R. 9510 Bill v4.0 (full bill) - Appendix L
%% Genre: hearing testimony + research-influence brief (opening statement, Q&A block,
%% research-influence matrix). Source deliverable 12. Table 11. Gray-scale Mermaid
%% figure (TikZ): Figure 7 (evidence to introduction handoff). Single hyphens only;
%% section symbol only in tables.
\billsec{APPENDIX L.}{TESTIMONY AND RESEARCH-INFLUENCE BRIEF (NON-OPERATIVE).}\phantomsection\label{toc:aL}

\uscprov{1}{\textit{Independent research draft. Not an enacted law, not pending
legislation, and not legal advice; not endorsed by the FDA, HHS, or any Member of
Congress. The bills and executive actions discussed are research influences only and
are not the legal basis of any operative provision.}}

\uscprov{1}{This appendix has two parts. Part I is hearing testimony, written to be
read aloud in about five minutes before the House Committee on Energy and Commerce
(Subcommittee on Health), with a question-and-answer block. Part II is a
research-influence brief that records, for the hearing record, exactly how recent
executive actions and emerging 119th Congress bills relate to H. R. 9510, and exactly
why none of them appears in its operative text.}

\uscprov{0}{\textbf{PART I. HEARING TESTIMONY.}}

\uscprov{1}{Chairman, Ranking Member, and Members of the Subcommittee: thank you for
the opportunity to testify on the verification of artificial intelligence in cancer
trials, and specifically on the safety of AI that does not merely advise a clinician
but physically moves around a patient and acts upon that patient.}

\uscprov{1}{\textbf{Who I am.} I am Kevin Kawchak, Chief Executive Officer of
ChemicalQDevice. Our work concerns the verification, validation, and uncertainty
quantification, or VVUQ, of high-consequence AI systems, including the control
software for surgical and patient-contacting robots used in oncology research. I appear
in an independent research capacity. The draft bill I will describe, H. R. 9510, is an
illustrative research instrument; no Member has agreed to sponsor it, and nothing I say
should be taken as the position of the FDA, of HHS, or of any Member of Congress.}

\uscprov{1}{\textbf{The problem.} Cancer trials are beginning to place robots in direct
physical contact with patients: surgical humanoids that cut and suture, positioning
systems that hold a patient during treatment, and platforms that place biopsy or
delivery needles. What such a robot does at the instant it touches a patient is decided
by software, and that software is increasingly written by AI in real time rather than
typed, line by line, by a human engineer. Here is the gap. No Federal law governs the
order of operations. A trial sponsor today may generate the robot-patient interaction
code, run it on a patient, and document the safety checks afterward, in whatever order
it chooses, because the order is simply not regulated. We call this the sequencing
gap. Sequence matters because of embodiment. In an ordinary data pipeline, a software
error produces a wrong number that a human can catch downstream before anyone is
harmed. In an embodied surgical system, the same class of error produces a wrong
physical motion that a human cannot catch in time, because one moving body concentrates
every failure mode into a single instrument acting in real time. Checking the code
after it has already moved is, in this setting, checking too late.}

\uscprov{1}{\textbf{The proposal.} H. R. 9510 supplies the one missing rule. It amends
the Federal Food, Drug, and Cosmetic Act and adds a single new section, 515D, codified
at 21 U.S.C. 360e-5, placed immediately after the predetermined change control
authority at section 515C that it most resembles. The core rule is one sentence of
policy: verify before generate. No robot-patient interaction code may be generated or
executed in a Physical AI oncology clinical investigation unless an automated VVUQ
process, bound to named external standards, has first cleared that interaction, with
the record documented and attested. A record created after generation or execution does
not satisfy the section. The timing of the record is itself the compliance fact. The
gate is concrete, not aspirational. Each check returns one of three results: ACCEPT,
BLOCK, or ESCALATE, which returns the interaction to a qualified human. For each
interaction the verification fraction must equal 1.0, meaning every check passes, with
no partial credit. Three gates carry a hard catastrophe predicate at perfect
agreement: no breach of a vascular no-fly volume, a guaranteed minimum clearance from
any human sharing the room, and a guaranteed safe state on fault or emergency stop
\cite{usc-360e4}.}

\uscprov{1}{\textbf{Why it is feasible.} This is not a hope; it has been run. In a
first study, VVUQ-01, an automated pipeline held the verification fraction at 1.0
across 51 of 51 tests, with 1 ACCEPT and 5 BLOCK outcomes. In a second study, VVUQ-02,
a ten-gate assurance suite over the control software for a surgical humanoid cleared
172 of 172 tests, with a composite mean of 93.56, reproducibly from a fixed seed,
20260525. The bill asks sponsors to do, as a precondition, what has already been
demonstrated to work. The gate is bound to recognized standards rather than to anyone's
discretion: its thresholds tie to ASME V\&V 40, IEC 80601-2-77, ISO/TS 15066, ISO
14971, NASA-STD-7009A, and IEEE 7009. That binding is what makes the test auditable and
what keeps it from becoming a moving target \cite{kawchak2026vvuq01, kawchak2026vvuq02}.}

\uscprov{1}{\textbf{Why it fits existing law.} I want to be candid about the policy
environment. The current Federal posture toward AI is deliberately deregulatory and
pro-innovation. This bill is designed to complement that posture, not to fight it. It
is not a broad AI mandate. It is a narrow, standards-bound safety gate for one thing:
embodied robot-patient code in clinical trials. The bill is additive. It leaves the
investigational-device exemption, human-subject protection, and the investigational new
drug framework in place. It frames the gate as a pre-authorized change-control
discipline of the kind section 515C already contemplates, routes its records through the
electronic-records rule and the quality-system regulation that sponsors already meet,
and adds a savings clause so existing State human-in-the-loop laws are preserved. It
creates no new agency and no new spending program; it directs rulemaking only
\cite{eo-14179, omb-m2521}.}

\uscprov{1}{\textbf{Respectful asks.} I respectfully ask the Subcommittee to consider
three things. First, treat the sequencing gap as a discrete, well-defined safety
question, separate from the larger and more contested debates about AI policy
generally. Second, recognize that verification before generation is technically
achievable today, on the evidence I have summarized, so the question before you is one
of timing and assurance, not of feasibility. Third, if the Subcommittee acts, keep any
measure narrow and standards-bound, so that it strengthens the credibility of American
oncology research without becoming the broad mandate that current Federal policy
disfavors. Thank you for your attention. I welcome your questions.}

\uscprov{1}{\textbf{Anticipated questions and answers.} The following exchange
anticipates the questions a Member is most likely to ask, with the responses the
witness is prepared to give.}

\uscprov{1}{\textit{Q. Is this a broad new mandate on artificial intelligence?}}

\uscprov{1}{A. No. The section reaches one thing only: robot-patient interaction code
in a Physical AI oncology clinical investigation. It does not touch disembodied
diagnostic software, clinical-decision-support tools that only advise a clinician, or
AI outside the trial setting. It is a single sequencing rule, bound to named external
standards, for code that physically moves around a patient.}

\uscprov{1}{\textit{Q. Will this slow down American innovation in cancer research?}}

\uscprov{1}{A. It is designed not to. The gate is automated, so it runs at machine
speed rather than waiting on manual review, and it asks sponsors to do as a
precondition what two studies already show is achievable. By making the verification
record auditable and standards-bound, it strengthens the credibility of the research it
governs, which is what allows results to be trusted and relied upon.}

\uscprov{1}{\textit{Q. Why legislate now rather than leave this to FDA guidance?}}

\uscprov{1}{A. The closest existing analogue, the predetermined change control plan
authority in section 515C and its final guidance, is voluntary, submission-based, and
neither automated nor required in advance of generation. Guidance can describe good
practice, but it cannot establish the order-of-operations rule that embodiment makes
necessary. The bill supplies that one rule and directs the Secretary to issue final
implementing regulations within 365 days \cite{usc-360e4}.}

\uscprov{1}{\textit{Q. Does the bill conflict with the current deregulatory posture or
with State law?}}

\uscprov{1}{A. It is built to complement the deregulatory posture: it is narrow,
standards-bound, and adds no new agency or spending program. On State law, it includes
a savings clause that preserves existing State human-in-the-loop and audit
requirements, so it sets a Federal floor for one safety question without displacing the
human-oversight protections States have already enacted \cite{eo-14365}.}

\uscprov{0}{\textbf{PART II. RESEARCH-INFLUENCE BRIEF.}}

\uscprov{1}{This brief records, for the hearing record, how each recent executive
action and each emerging 119th Congress bill relates to H. R. 9510, and demonstrates
that none of them is, or could be, the legal basis of any operative provision. Table 11
is the research-influence matrix; the placement column states where each authority is
used, and the operative-versus-influence rule keeps every one of them out of the
statutory text and the comparative print.}

{\footnotesize
\begin{xltabular}{\textwidth}{@{}L{2.0cm} L{3.4cm} Y L{2.4cm}@{}}
\hline\noalign{\vskip 0.04cm}
\textbf{Authority} & \textbf{What it is} & \textbf{What it informs} & \textbf{Placement}\\
\hline\noalign{\vskip 0.04cm}
\endfirsthead
\hline\noalign{\vskip 0.04cm}
\textbf{Authority} & \textbf{What it is} & \textbf{What it informs} & \textbf{Placement}\\
\hline\noalign{\vskip 0.04cm}
\endhead
\hline\noalign{\vskip 0.04cm}
\endlastfoot
EO 14179 & Removing Barriers to American Leadership in AI (90 FR 8741) & Deregulatory
posture in the findings narrative; the gate framed as credibility rather than mandate &
Memo; matrix; testimony\\
\hline\noalign{\vskip 0.04cm}
EO 14365 & Ensuring a National Policy Framework for AI & State-law review and preemption
context; the case for one national documentation standard & Appendix; matrix;
testimony\\
\hline\noalign{\vskip 0.04cm}
OMB M-25-21 & Accelerating Federal Use of AI (pre-deployment testing) & The
pre-deployment-testing analogue; the Federal lifecycle the gate maps onto & Memo;
matrix; testimony\\
\hline\noalign{\vskip 0.04cm}
HHS AI Strategy & The HHS Artificial Intelligence Strategy (Dec. 2025) & Governance and
reproducibility framing that resonates with VVUQ goals & Memo; matrix\\
\hline\noalign{\vskip 0.04cm}
H.R. 238 & Healthy Technology Act of 2025 & AI prescribing-authority direction; why
autonomous clinical AI raises the evidentiary bar & Testimony; matrix\\
\hline\noalign{\vskip 0.04cm}
S. 1399 & Health Tech Investment Act & Medicare payment-class direction;
clearance-gated payment makes validation the entry ticket & Testimony; matrix\\
\hline\noalign{\vskip 0.04cm}
H.R. 6361 & Ban AI Denials in Medicare Act & The human-over-AI signal; bipartisan
distrust of opaque AI decisions & Findings narrative; matrix\\
\hline\noalign{\vskip 0.04cm}
H.R. 5045 & HEALTH AI Act & Generative-AI research direction; potential funding for
validation methodology & Memo; matrix\\
\hline\noalign{\vskip 0.04cm}
H.Res. 694 & Sense of the House on the CMS WISeR AI pilot & The countervailing
human-review signal; demand for auditable AI & Findings narrative; matrix\\
\hline
\end{xltabular}}
\vspace{-0.8cm}
\figcaption{Table 11. The research-influence matrix: emerging 119th Congress bills and
executive actions, what each informs, and where each is used. The placement column is
memo, appendix, or testimony only; none is ever in the operative text.}

\uscprov{1}{\textbf{How each influence relates, and how it stays out of operative
text.} The executive actions set the environment the bill is drafted into. EO 14179 and
EO 14365 together describe a deregulatory, preemption-minded Federal posture; they
influence the bill in one direction only, arguing against a broad mandate and in favor
of a narrow, voluntary-credibility framing and a single national documentation
standard. OMB M-25-21 is more structural: its requirement of pre-deployment testing, an
impact assessment, and ongoing monitoring is the cleanest surviving Federal
articulation of a VVUQ-style lifecycle, used as an analogue to show that verifying
before deployment is already a recognized Federal idea. The HHS Artificial Intelligence
Strategy of December 2025 is invoked only to show alignment of values. None of these is
operative authority, and the bill never treats them as such. They are executive and
sub-regulatory instruments, not statutes; they can change with an administration; and a
safety gate that drew its legal force from an executive order would be unstable by
construction. The bill therefore takes its operative force from the Federal Food, Drug,
and Cosmetic Act and from recognized consensus standards
\cite{eo-14179, eo-14365, omb-m2521, hhs-ai-strategy-2025}.}

\uscprov{1}{The emerging bills establish legislative momentum on both sides of one
question. On the autonomy side, H.R. 238 would let AI qualify as a practitioner
eligible to prescribe, and S. 1399 would create a Medicare new-technology payment class
for cleared algorithm-based services; both raise the stakes of getting autonomous
clinical AI right, which is the gap a verification gate addresses. On the human-
oversight side, H.R. 6361 would bar AI-driven prior authorization in traditional
Medicare, and H.Res. 694 urges CMS to halt an AI coverage pilot; both reflect a
bipartisan insistence that a human stay accountable for consequential AI decisions.
H.R. 5045 sits alongside these as a research-funding vehicle that could support
validation and evaluation methodology. These bills inform the testimony, the findings,
and the matrix, and nothing more. They are introduced or in-committee measures; not one
is enacted, and an introduced bill has no legal effect, so citing one as the basis of
an operative clause would be both incorrect and misleading
\cite{hr238-healthy-tech, s1399-health-tech-investment, hr6361-ban-ai-denials, hr5045-health-ai, hres694-wiser}.}

\uscprov{1}{\textbf{Verification note.} Each authority above was verified by number,
subject, and status before use, and each is presented strictly as a research influence;
readers should re-verify currency before relying on any entry. Three cautions govern
this brief. First, there is no valid Act in the 119th Congress on this subject; the
measures titled VALID Act of 2025 are an unrelated matter, so no such Act is cited
here. Second, a standalone AI-in-clinical-trials bill and a standalone Federal
health-AI liability bill could not be verified, and no bill numbers are assigned to
them. Third, the sponsor of H.Res. 694 was not cleanly confirmable, so it is cited by
number and subject only. Consistent with the whole package, only enacted statutes,
in-force rules, and recognized standards carry operative weight; the bills and
executive actions in this brief carry none.}

\uscprov{1}{\textbf{Evidence to introduction.} Figure 7 traces the handoff this
appendix sits within: the author sets requirements; Claude Code Opus 4.8 generates and
runs the work; the VVUQ gates and tests confirm it; the Office of the Legislative
Counsel puts the result in legislative form; and a sponsoring Member introduces it
through the House eHopper \cite{anthropic-claude-code, house-ehopper}.}

\begin{mermaidfig}
\node[mmperson,text width=20mm] (k) at (0,0) {Author (human)};
\node[mmstep,text width=20mm] (ai) at (2.8,0) {Claude Code Opus 4.8};
\node[mmdec,text width=20mm] (t) at (5.6,0) {VVUQ gates and tests};
\node[mmstep,text width=20mm] (l) at (8.4,0) {Legislative Counsel};
\node[mmgoal,text width=20mm] (h) at (11.2,0) {House (eHopper)};
\draw[mmedge] (k) -- node[mmlabel]{1 requirements} (ai);
\draw[mmedge] (ai) -- node[mmlabel]{2 generate, run} (t);
\draw[mmedge] (t) -- node[mmlabel]{3 bill, deliverables} (l);
\draw[mmedge] (l) -- node[mmlabel]{4 introduce} (h);
\end{mermaidfig}
\vspace{-0.3cm}
\figcaption{Figure 7. Evidence to introduction handoff: requirements to autonomous
generation and gating, to legislative form, to introduction by a sponsoring Member.}
